\documentclass[11pt]{article}
\usepackage[margin=1in]{geometry}
\usepackage{amsmath,amssymb,amsthm,hyperref}
\newtheorem{theorem}{Theorem}
\title{Chambolle--Pock convergence with step-size product four}
\author{ziangni-sys}
\date{Conjecture 00000008833}
\begin{document}
\sloppy
\maketitle
\noindent Both originals assert that Chambolle--Pock convergence always
requires the product of the two step sizes to be at most one. No
normalization of the coupling operator is imposed. The following actual
scalar convex problem disproves this necessary-condition claim.

\section*{Problem and exact algorithm}
Use the real Hilbert spaces $X=Y=\mathbb R$, and set
\[
 f(x)=\tfrac12x^2,\qquad g(z)=\tfrac12z^2,\qquad Kx=\tfrac14x,
 \qquad \tau=\sigma=2,\qquad\theta=1.
\]
Both objectives are finite, continuous, proper and convex. The coupling is
nonzero and bounded, with $K^*=K$. Completing the square gives
\[
 g^*(y)=\sup_z\{yz-\tfrac12z^2\}=\tfrac12y^2,
 \qquad yz-\tfrac12z^2=\tfrac12y^2-\tfrac12(z-y)^2.
\]
The supremum is attained at $z=y$. For either proximal objective at step
size two, completing the square gives the unique minimizer
\[
 \mathop{\rm prox}_{2f}(a)=\mathop{\rm prox}_{2g^*}(a)=a/3,
 \qquad \tfrac12u^2+\tfrac14(u-a)^2
       =\tfrac34(u-a/3)^2+\tfrac16a^2.
\]
The saddle function is $L(x,y)=x^2/2+xy/4-y^2/2$.
The point $(0,0)$ is a saddle because $L(0,y)\le L(0,0)\le L(x,0)$.

Algorithm 1 of Chambolle--Pock~\cite{CP} updates
\[
 y'=\mathop{\rm prox}_{\sigma g^*}(y+\sigma K\bar x),\quad
 x'=\mathop{\rm prox}_{\tau f}(x-\tau K^*y'),\quad
 \bar x'=x'+\theta(x'-x).
\]
Consequently this instance has the exact update
\[
 y'=(y+\bar x/2)/3,\qquad x'=(x-y'/2)/3,
 \qquad \bar x'=2x'-x=-x/3-y'/3.
\]
The usual initialization $\bar x^0=x^0$ is included below.

\begin{theorem}
Every initial state $(x^0,y^0,\bar x^0)$ converges to $(0,0,0)$ under
this iteration, although $\tau\sigma=4>1$.
\end{theorem}
\begin{proof}
Let $M=\max(|x|,|y|,|\bar x|)$. The update satisfies
\[
 |y'|\le M/2,\qquad |x'|\le 5M/12\le M/2,
 \qquad |\bar x'|\le M/3+M/6=M/2.
\]
Thus $M_n\le 2^{-n}M_0\to0$. Every coordinate converges to zero,
so the primal-dual pair converges to the exhibited saddle from every
initial pair. The positive step sizes have product four, contradicting
the claimed universal necessity of a product at most one.
\end{proof}

\newpage
\section*{Scope and the standard operator-norm condition}
The counterexample concerns exactly the unqualified numerical product
bound stated in both source versions. The operator $K$ is nonzero, the
objectives meet the standard convex assumptions, and a saddle exists.
Convergence is proved for every initial state, rather than only a fixed
point initialization or a specially selected orbit. The extrapolation
parameter is the standard value $\theta=1$.

The usual sufficient condition in Chambolle--Pock's theorem is
$\tau\sigma\|K\|^2<1$~\cite{CP}. Here $\|K\|=1/4$, so that product
is $1/4$. Our result is therefore consistent with the standard theorem.
The source omits this operator-norm factor and makes its step-product
condition universal; it does not assume $\|K\|=1$. The direct proof above
does not invoke the convergence theorem. Failure of this first conjunct
suffices to disprove the conjunction; the separate preconditioning claim
need not be resolved.

\section*{Formal objects and complete coverage}
The Lean file defines the actual real quadratic \texttt{energy} and
proves \texttt{Continuous} and \texttt{ConvexOn} properties.
Since it is real-valued everywhere, its properness has no exceptional
infinite values. The operator \texttt{coupling} is a genuine bounded
\texttt{ContinuousLinearMap} on the real Hilbert line. Its scalar action,
nonzeroness and adjoint identity against every pair of vectors are proved.

\texttt{fenchel\_conjugate} computes the genuine \texttt{sSup} of the
range $z\mapsto yz-f(z)$. A universal upper bound and an attaining point
establish boundedness, existence and the value of this supremum.
\texttt{prox\_exact} proves an if-and-only-if characterization of all
global minimizers of the actual proximal objective. Thus the proximal
maps used by the algorithm are derived, including their uniqueness.
\texttt{zero\_saddle} checks the saddle inequalities for all $x,y$.

\texttt{step} implements both proximal updates, the bounded coupling
operator and the extrapolation step with the stated parameters. The
adjoint identity justifies using the same scalar coupling in the primal
update. \texttt{step\_formula} derives the displayed recurrence.
The orbit is defined by actual function iteration, and \texttt{size} is
the maximum of all three absolute coordinate values.
\texttt{one\_step} and \texttt{size\_contracts} prove its contraction.
\texttt{orbit\_bound} proves the bound $2^{-n}M_0$ for every natural $n$.
Finally, \texttt{orbit\_converges} proves convergence in the ordinary
product topology for every initial triple, and
\texttt{product\_bound\_counterexample} combines that convergence with
positive steps, nonzero coupling and $\tau\sigma>1$.

\section*{Reproduction}
Lean 4.19.0 and Mathlib commit
\begin{center}\small\texttt{c44e0c8ee63ca166450922a373c7409c5d26b00b}\end{center}
are pinned. From \texttt{lean/}, run \texttt{lake update},
\texttt{lake exe cache get}, \texttt{lake build}, and
\texttt{lake env lean Main.lean -DwarningAsError=true}.
The six printed dependency audits use only standard logical axioms.
No admitted statements or auxiliary computations are used.

\begin{thebibliography}{1}
\bibitem{CP} A. Chambolle and T. Pock,
\emph{A first-order primal-dual algorithm for convex problems with applications to imaging},
J. Math. Imaging Vis. 40 (2011), 120--145. Algorithm 1 and Theorem 1.
\href{https://www.cmap.polytechnique.fr/preprint/repository/685.pdf}{Author preprint}.
\end{thebibliography}
\end{document}
